\subsubsection{Bounded timetable feasibility and feature summaries}
\label{sec:reduced-od-raptor}

RAPTOR denotes \emph{Round-Based Public Transit Routing}.  It is a
timetable-routing method for finding feasible public-transport journeys and
their earliest arrival times without first converting the timetable into a
time-expanded graph.  Its state is organized by stop and round: round zero
contains the origin and its walking reachability, and round $r$ contains labels
reachable with at most $r$ vehicle boardings.  Routes reachable from labels
improved in the preceding round are scanned in timetable order; newly improved
stop labels are then propagated through transfer footpaths.  Keeping the best
arrival label for each transfer count yields an arrival-versus-transfers Pareto
frontier.  Unlike graph-based shortest-path algorithms, the standard method
operates directly on route and timetable arrays and does not require a global
priority queue over timetable events.

In this project RAPTOR is a preprocessing algorithm, not an OD estimator.  Its
role is to decide which destinations are reachable from an origin and departure
time and to summarize the scheduled journeys used later to construct choice
sets and measurement responses.  Passenger counts and demand parameters do not
enter the routing calculation.

The network-level preprocessing therefore begins with a deterministic
RAPTOR-style search implemented in
\path{public_transportation.preprocessing.reduced_od.raptor}.  It consumes the
compact timetable arrays described above, explicit directed footpaths, a
physical origin stop, a departure time, and a maximum number of transfers.  It
does not construct a global time-expanded graph and does not call either
assignment backend.

A search round represents exactly one vehicle boarding.  With maximum transfer
count $K$, the algorithm executes $K+1$ route rounds.  Each scheduled trip is
scanned once per round in stop order.  The scan retains its best currently
boardable label and propagates it to downstream stops, so its work is linear in
the stop-time records scanned rather than quadratic in possible board/alight
pairs.  Positive directed footpaths are relaxed without consuming a vehicle
round.  Transfer depth is consequently bounded by construction.

Within a round, the retained label is the deterministic earliest-arrival label,
with waiting time, walking time, and path identity used as tie breakers.  Across
rounds, the result retains the Pareto frontier in arrival time and number of
boardings.  Every label reports total elapsed time, waiting, walking,
in-vehicle time, transfers, and its scheduled transit legs.  The feature
identity
\[
T_{\mathrm{elapsed}}
=T_{\mathrm{wait}}+T_{\mathrm{walk}}+T_{\mathrm{vehicle}}
\]
therefore remains directly auditable.  Phase~4 does not enumerate every path
with a distinct waiting/walking tradeoff inside one transfer round; bounded
choice-set enumeration and pruning belong to Phase~5.  The current code is
described as RAPTOR-style because it scans every scheduled trip in every bounded
round and does not yet apply the standard marked-stop/marked-route pruning or a
route-pattern search for the earliest boardable trip.  These omissions affect
preprocessing speed, not the transfer-round semantics or the returned
earliest-arrival labels.

For every absent destination, a topology-only zero/one shortest-path pass over
route reachability and footpaths records one fail-closed reason:
\texttt{no\_topological\_\allowbreak path},
\texttt{exceeds\_transfer\_\allowbreak limit}, or
\texttt{no\_timetable\_\allowbreak feasible\_\allowbreak journey}.  The origin itself is labeled
separately and is not presented as a destination.  This distinction prevents a
missed scheduled connection from being misreported as a disconnected network.
Independent range queries sort and deduplicate their departure instants and
return one complete result per instant.

The Phase-4 CPU benchmark used overlapping patterns of at most 20 stops, four
departures per pattern, and a three-transfer bound.  For 25, 50, 100, and 200
physical stops, mean query times were approximately 0.00065, 0.00091, 0.00124,
and 0.00187 seconds, corresponding to about 1,535, 1,101, 803, and 534 queries
per second.  The compact timetable arrays occupied 3,656, 7,208, 14,664, and
29,576 bytes.  These machine-dependent measurements are diagnostic rather than
test thresholds; full-network performance must still be measured on its actual
pattern, trip, period, and query counts.

